\documentclass[10pt,a4paper]{amsart}
\usepackage[margin=19mm]{geometry}
\usepackage{amsmath,amssymb,mathtools}
\usepackage[scr=rsfs,cal=euler]{mathalfa}
\usepackage{xfrac}
\usepackage[unicode,hidelinks,bookmarks=false]{hyperref}
\newcommand{\ie}{i.e.\ }
\newcommand{\cf}{cf.\ }
\newcommand{\resp}{respectively\ }
\newcommand{\Subsection}{\subsection}
\providecommand{\nobreakdash}{\nobreak\hbox{-}}
\setlength{\emergencystretch}{2em}
\multlinegap0pt
\usepackage{bm}
\usepackage{bbm}
\usepackage{dsfont}
\usepackage{xspace}
\usepackage{cancel}
\usepackage{mathtools}
\usepackage{amsthm}
\makeatletter
\@ifundefined{theorem}{\newtheorem{theorem}{Theorem}}{}
\@ifundefined{lemma}{\newtheorem{lemma}[theorem]{Lemma}}{}
\makeatother
\DeclareMathOperator{\Tr}{Tr}
\newcommand{\bim}{\kappa}
\newcommand{\mattyone}{X_1}
\newcommand{\mattytwo}{X_2}
\newcommand{\simpkill}{\kappa}
\title[Conjugate points on Lie groups with left-invariant metrics]{Conjugate points on Lie groups with left-invariant metrics}
\author{Alice Le Brigant, Leandro Lichtenfelz, Stephen C. Preston}
\date{Source: arXiv:2408.03854v3 (CC BY 4.0)}
\begin{document}
\maketitle

\noindent\emph{Source and licence.} Conjugate points on Lie groups with left-invariant metrics. Version arXiv:2408.03854v3 (CC BY 4.0), distributed under the Creative Commons Attribution 4.0 International licence, as recorded in the arXiv licence metadata for this exact version and shown on the abstract page https://arxiv.org/abs/2408.03854v3. Full rights notice: licenses/arxiv-2408.03854-CC-BY-4.0.txt.\par\smallskip

\noindent\textit{Editorial scope.} Counted unit: the Lemma whose source label is \texttt{sl2conditionsecondtime} in the article (source0.tex lines 1090--1096, source label \texttt{sl2conditionsecondtime}), delivered together with the article's own complete original proof of it (source0.tex lines 1100--1173, one \texttt{proof} environment of 74 lines). No mathematics is written, repaired or re-derived: statement and proof are the article's own text line for line. The proof is detached from the statement in the source: the article states the result at lines 1090--1096 and prints its proof at lines 1100--1173, and the proof environment's own header names the statement, so the association is the article's own. Required context retained uncounted: none. Every definition, display and label of those lines is retained unchanged. Omitted from this package and not redistributed: 1--1089, 1097--1099, 1174--1303. Nothing inside a retained excerpt is omitted or altered: every retained body is delivered exactly as the article's lines are written, so no omission receipt of any kind accompanies this package. Citations: the retained excerpts carry no citation command at all, so no bibliography entry of the article is transcribed and no reference is redistributed. Reference adaptation: none was needed and none was made; every reference inside the delivered unit resolves inside it, so no unresolved reference or citation is produced. Printed slips kept literal: no printed word, symbol, hypothesis or number of the counted statement or of its proof was altered. Layout and class: the article's own class is \texttt{amsart} with packages including amsmath, amsfonts, amssymb, mathrsfs, amsthm, graphicx, xcolor, array, comment, geometry, hyperref, float, caption, cancel; the wrapper is the lane's pinned \texttt{amsart} editorial wrapper at 10pt on A4, which restates the theorem-like environments the retained text needs and adds the article's own macro definitions that the retained text needs (\texttt{\textbackslash Tr}, \texttt{\textbackslash bim}, \texttt{\textbackslash mattyone}, \texttt{\textbackslash mattytwo}, \texttt{\textbackslash simpkill}), with extra wrapper packages (bm, bbm, dsfont, xspace, cancel, mathtools). The delivered numbering is the unit's own. Assets: no retained span contains a figure environment, a graphics-inclusion command, a PDF-page inclusion, a table or a TikZ picture; no image file is copied into this package, the package declares no asset dependency and no image file is redistributed with it. Licence: version arXiv:2408.03854v3 of this article is distributed under the Creative Commons Attribution 4.0 International licence, as recorded in the arXiv licence metadata for this exact version and shown on the abstract page https://arxiv.org/abs/2408.03854v3. A full rights notice is delivered at licenses/arxiv-2408.03854-CC-BY-4.0.txt. Characters: every retained excerpt is pure ASCII, so the delivered engine, XeLaTeX with its default Latin Modern text font, reports no missing character.
\begin{lemma} 
    Let $\simpkill(\mattyone,\mattytwo)=\tfrac{1}{2}\Tr(\mattyone\mattytwo)$ denote the Killing form on $\mathfrak{sl}_2(\mathbb{R})$. Then for any distinct matrices $\mattyone,\mattytwo\in\mathfrak{sl}_2(\mathbb{R})$, the equation $$\det{(e^{t\mattyone}-e^{t\mattytwo})} = 0$$
    has a solution for $t>0$ if and only if at least one of $\simpkill(\mattyone,\mattyone)$ or $\simpkill(\mattytwo,\mattytwo)$ is negative or, if both are nonnegative, we have
\begin{equation}\label{sl2conditionsecondtime}
2\sqrt{\simpkill(\mattyone,\mattyone)\simpkill(\mattytwo,\mattytwo)} < 2\simpkill(\mattyone,\mattytwo) < \simpkill(\mattyone,\mattyone)+\simpkill(\mattytwo,\mattytwo).
    \end{equation}
\end{lemma}

\begin{proof}
    Since $\Tr{\mattyone}=0$ we have by the Cayley-Hamilton theorem that $\mattyone^2 = -(\det{\mattyone}) I$ so that $\simpkill(\mattyone,\mattyone) = -\det{\mattyone}$, and similarly for $\mattytwo$. Let 
    $$d(t) := \det{(e^{t\mattyone}-e^{t\mattytwo})}.$$
 

\textbf{Case 1: $\simpkill(\mattyone,\mattyone)=a^2$, $\simpkill(\mattytwo,\mattytwo)=b^2$, $a,b>0$.} 

    In this case we must have $a\ne b$, for otherwise the strict inequality condition \eqref{sl2conditionsecondtime} cannot be satisfied.
    The matrix exponentials are given by 
    $$ e^{t\mattyone} = (\cosh{at}) I + a^{-1}(\sinh{at}) \mattyone, \qquad e^{t\mattytwo} = (\cosh{bt}) I + b^{-1}(\sinh{bt}) \mattytwo. $$
    It is then straightforward to compute that 
\begin{equation*}%\label{determinantdifference}
    d(t) = 2\big( 1-\cosh{at}\cosh{bt} + D \sinh{at}\sinh{bt}\big), \qquad D:=\simpkill(\mattyone,\mattytwo)/(ab), 
    \end{equation*}  
If $D\le 1$, then for all $t>0$, we have 
    $$ d(t) \le 2(1-\cosh{at}\cosh{bt} + \sinh{at}\sinh{bt}) = 2\big(1 - 
    \cosh{\big((a-b)t\big)}\big) < 0. $$
    If $D\ge \tfrac{1}{2}(a^2+b^2)$ then we can compute that \begin{align*}
        d''(t) &= (2Dab-a^2-b^2) \cosh{at}\cosh{bt} + \big( (a^2+b^2)D - 2ab\big) \sinh{at}\sinh{bt} \\
        &\ge \tfrac{1}{2} (a^2-b^2)^2 \sinh{at}\sinh{bt}\ge 0,
        \end{align*}
    so $d(t)$ is positive for all $t>0$. 
    On the other hand if $1<D<\tfrac{1}{2ab}(a^2+b^2)$, then $d(0)=d'(0)=0$ while $$d''(0) = 2(2abD-a^2-b^2) < 0,$$
    so that $d(t)$ is negative for small positive $t$,
    while for large $t$ we have $d(t) \approx \tfrac{1}{2}(D-1)e^{(a+b)t}>0$. Hence there must be a change of sign for some positive $t$.

    \textbf{Case 2: $\simpkill(\mattyone,\mattyone)=-a^2$, $\simpkill(\mattytwo,\mattytwo)=b^2$, $a,b>0$.} 
    
    Then we have 
    $$d(t) = 2(1-\cos{at} \cosh{bt} + D\sin{at} \sinh{bt}), \qquad D = \simpkill(\mattyone,\mattytwo)/(ab). $$
    For $\tau:=\pi/a$, we clearly have $d(\tau)>0$ and $d(2\tau)<0$, so we get a crossing for any values of $a$ and $b$. 

    \textbf{Case 3: $\simpkill(\mattyone,\mattyone)=-a^2$, $\simpkill(\mattytwo,\mattytwo)=-b^2$, $a,b>0$.}  

    We can write
    \begin{align*}
        d(t) &= 2-2\cos{a t}\cos{bt} +2 D \sin{at}\sin{bt}, \qquad D:=\simpkill(\mattyone,\mattytwo)/(ab) \\
        &= 
        2(1+D) \sin^2{\left( \frac{a+ b}{2} \, t\right)} + 2(1-D)\sin^2{\left( \frac{ b -a}{2} \, t\right)}.
    \end{align*}
    Note that $\lvert D\rvert\ge 1$ in this case with equality if and only if $\mattytwo$ is a multiple of $\mattyone$. To see this, apply conjugation to get a canonical form for $\mattyone$, using invariance of the Killing form, to assume without loss of generality that
    $$ \mattyone = \left( \begin{matrix} 0 & r \\ 
    -r & 0 \end{matrix}\right), \qquad \mattytwo = \left(\begin{matrix} x & y+z \\ y-z & -x\end{matrix}\right)$$
    for some reals $r,x,y,z$. 
    Then we have that 
    $$ r^2=a^2, \qquad \simpkill(\mattyone,\mattytwo) = rz, \qquad 
    x^2+y^2-z^2= -b^2.$$
    Hence $$\lvert D\rvert = \lvert rz/ab\rvert = \lvert z/b\rvert \ge 1,$$
    with equality iff $x=y=0$, in which case $\mattytwo$ is a multiple of $\mattyone$.
If $a=b$ then $\lvert D\rvert > 1$ since $\mattyone\ne \mattytwo$, and we have $d(t) = 2(1+D) \sin^2{at}$ which obviously vanishes at all integer multiples of $\pi/a$. 
Otherwise we may assume without loss of generality that $b>a$, and  defining $\tau_1 = \frac{2\pi}{b+a}$ and $\tau_2 = \frac{2\pi}{b-a}$, we will have 
    $$d(\tau_1) = 2(1-D)  \qquad \text{and}\qquad d(\tau_2) = 2(1+D).$$
    Since $\lvert D\rvert > 1$, these values must have opposite signs,
    so $d(t)$ vanishes between $\tau_1$ and $\tau_2$. 
 

\textbf{Case 4: $\simpkill(\mattyone,\mattyone)=0$, $\simpkill(\mattytwo,\mattytwo)=b^2$, $b>0$.}

In this case we have 
$$ d(t) = 2(1-\cosh{bt} + Dt \sinh{bt}), \qquad D:= \bim(\mattyone,\mattytwo)/b.$$
The same reasoning as in Case 1 shows that this can vanish if and only if $0<2D<b$, which is precisely condition \eqref{sl2conditionsecondtime}.

\textbf{Case 5: $\simpkill(\mattyone,\mattyone)=0$, $\simpkill(\mattytwo,\mattytwo)=-b^2$, $b>0$.}

In this case we have 
$$d(t) = 2(1-\cos{bt} + 2Dt\sin{bt}) = 4\sin{(\tfrac{bt}{2})} \big( Dt \cos{(\tfrac{bt}{2})} + \sin{(\tfrac{bt}{2})}\big), \qquad D:= \bim(\mattyone,\mattytwo)/b.$$
This obviously vanishes at all integer multiples of $2\pi/b$. 

\textbf{Case 6: $\simpkill(\mattyone,\mattyone)=0$, $\simpkill(\mattytwo,\mattytwo)=0$.}

In this last case we have 
$$d(t) = 2t^2 \bim(\mattyone,\mattytwo),$$
which obviously never vanishes; also the strict inequality of condition \eqref{sl2conditionsecondtime} cannot be satisfied.
\end{proof}

\end{document}
